\documentclass[11pt,a4paper]{article}
\usepackage[utf8]{inputenc}
\usepackage[T1]{fontenc}
\usepackage{lmodern}
\usepackage{amsmath,amssymb,amsfonts}
\usepackage{graphicx}
\usepackage{geometry}
\usepackage{booktabs}
\usepackage{hyperref}
\usepackage{float}
\usepackage{subcaption}
\usepackage{xcolor}
\usepackage{listings}
\usepackage{enumitem}
\usepackage{fancyhdr}
\usepackage{titlesec}

\geometry{margin=1in}
\hypersetup{colorlinks=true, linkcolor=blue, citecolor=blue, urlcolor=blue}

\pagestyle{fancy}
\fancyhf{}
\fancyhead[L]{\nouppercase{\leftmark}}
\fancyhead[R]{\thepage}
\renewcommand{\headrulewidth}{0.4pt}

\title{Magic Squares: Geometric and Manifold Analysis}
\author{Magic Squares Research Team}
\date{March 04, 2025}

\begin{document}

\maketitle
\tableofcontents
\newpage
\section{Introduction}

This report presents a comprehensive analysis of magic squares, focusing on their geometric properties, 
Earth Mover's Distance (EMD) relationships, and manifold structure. The analysis pipeline includes 
multiple stages of data processing and visualization to extract meaningful patterns and insights.

\begin{figure}[H]
    \centering
    \includegraphics[width=0.9\textwidth]{figures/historical_context.png}
    \caption{Historical context of magic squares, showing the Lo Shu square (c. 190 BCE), 
    Dürer's magic square from Melencolia I (1514), and a timeline of major developments.}
    \label{fig:historical_context}
\end{figure}

\begin{figure}[H]
    \centering
    \includegraphics[width=0.9\textwidth]{figures/research_framework.png}
    \caption{Research framework showing the analysis pipeline for magic squares, 
    including geometric analysis, EMD analysis, manifold learning, and predictive modeling.}
    \label{fig:research_framework}
\end{figure}

\section{Methodology}

Our analysis follows a multi-stage pipeline that examines magic squares from different perspectives:

\begin{enumerate}
    \item \textbf{Geometric Analysis:} We represent magic squares as point clouds in 3D space and analyze 
    their geometric properties, including convex hull characteristics, point distributions, and eigenvalue patterns.
    \item \textbf{Earth Mover's Distance (EMD) Analysis:} We compute the optimal transport cost between magic squares 
    to quantify their similarity and analyze the resulting distance matrix.
    \item \textbf{Manifold Learning:} We apply dimensionality reduction techniques to explore the underlying 
    structure of the magic square space and estimate its intrinsic dimension.
    \item \textbf{Predictive Modeling:} We build machine learning models to predict geometric properties 
    and identify the most important features that characterize magic squares.
\end{enumerate}

\section{Geometric Analysis}

We analyze the geometric properties of magic squares by representing them as point clouds in 3D space. 
Each point corresponds to a cell in the magic square, with its position determined by the cell's coordinates 
and value. We then compute various geometric measures, including convex hull volume, surface area, and point distributions.

\section{Earth Mover's Distance Analysis}

Earth Mover's Distance (EMD) provides a natural way to measure the similarity between magic squares 
by computing the optimal transport cost between them. We analyze the pairwise EMD distances between 
magic squares and explore the resulting distance matrix to identify clusters and patterns.

\section{Manifold Analysis}

We explore the manifold structure of the magic square space using dimensionality reduction techniques 
such as t-SNE. We estimate the intrinsic dimension of the manifold and analyze its local neighborhood 
structure to gain insights into the underlying geometry of magic squares.


## Manifold Analysis

### Dimension Estimates


### Local Structure


### Curvature Analysis


\begin{figure}[H]
    \centering
    \includegraphics[width=0.8\textwidth]{figures/manifold_structure.png}
    \caption{Analysis of the manifold structure of magic squares, including intrinsic dimension estimates, 
    local neighborhood structure, and persistence diagrams.}
    \label{fig:manifold_structure}
\end{figure}

\section{Predictive Modeling}

We build machine learning models to predict geometric properties of magic squares based on their structural features. 
We use random forests, gradient boosting, and Lasso regression to identify the most important features and evaluate 
their predictive performance.


## Predictive Modeling Results

Our predictive modeling analysis reveals several key findings:

1. Best Performance: The model achieved the highest accuracy in predicting volume, with an R² score of 0.924 and mean squared error of 5.998.

2. Model Performance by Target:
   - volume: R² = 0.924, MSE = 5.998
   - area: R² = 0.864, MSE = 8.790
   - mean_dist: R² = 0.364, MSE = 0.001
   - std_dist: R² = 0.356, MSE = 0.003

The moderate predictive performance indicates consistent but complex relationships between geometric features and target properties.

\begin{figure}[H]
    \centering
    \includegraphics[width=0.8\textwidth]{figures/feature_importance_0.png}
    \caption{Feature importance for predicting geometric properties of magic squares, 
    showing the relative importance of different structural features.}
    \label{fig:feature_importance_0}
\end{figure}

\begin{figure}[H]
    \centering
    \includegraphics[width=0.8\textwidth]{figures/feature_importance_1.png}
    \caption{Feature importance for predicting geometric properties of magic squares, 
    showing the relative importance of different structural features.}
    \label{fig:feature_importance_1}
\end{figure}

\begin{figure}[H]
    \centering
    \includegraphics[width=0.8\textwidth]{figures/feature_importance_2.png}
    \caption{Feature importance for predicting geometric properties of magic squares, 
    showing the relative importance of different structural features.}
    \label{fig:feature_importance_2}
\end{figure}

\begin{figure}[H]
    \centering
    \includegraphics[width=0.8\textwidth]{figures/feature_importance_3.png}
    \caption{Feature importance for predicting geometric properties of magic squares, 
    showing the relative importance of different structural features.}
    \label{fig:feature_importance_3}
\end{figure}

\section{Conclusion}

Our comprehensive analysis of magic squares reveals interesting patterns and relationships in their geometric properties, 
EMD distances, and manifold structure. The predictive modeling results highlight the most important features that 
characterize magic squares and provide insights into their underlying structure.

Future work could explore more advanced geometric measures, apply different manifold learning techniques, 
and investigate the relationship between magic squares and other mathematical objects.

\end{document}